% ============================================================================
% Fichier  : partie3/P03_C11_privileges.tex
% Titre    : Élévation de Privilèges
% Auteur   : MINKA MI NGUIDJOÏ Thierry Emmanuel
% Version  : 1.0 -- novembre 2026
% ============================================================================
% Sources : MITRE ATT&CK (tactique Privilege Escalation) ; CIS Benchmarks ;
% documentation de sudo et des capacités Linux. Opération LOGONE : cas fictif.
% ============================================================================

\chapter{Élévation de Privilèges}
\label{chap:privileges}
\chaptermark{Élévation de privilèges}

\epigraph{%
    « L'élévation de privilèges exploite rarement le noyau~;\\
    elle exploite le plus souvent une négligence de configuration. »%
}{Principe du cours}

\niveauChapitre{Intermédiaire à expert}{Chapitre~\ref{chap:exploitation}~;
cours SEC3266, chapitre sur les systèmes.}

\begin{boxobjectifs}
\begin{itemize}
    \item Distinguer élévation verticale et mouvement horizontal.
    \item Recenser les vecteurs d'élévation sous Linux et sous Windows.
    \item Lire un relevé de configuration pour y repérer un vecteur.
    \item Classer les vecteurs par probabilité et par risque de dommage.
    \item Proposer la correction et la détection de chaque vecteur.
\end{itemize}
\end{boxobjectifs}

\bigskip

% ============================================================================
\section*{Question directrice : un compte web limité peut-il devenir administrateur du serveur~?}
% ============================================================================

Après la démonstration sur l'ancien portail (Chapitre~\ref{chap:exploitation}),
M.~TAMO dispose d'un compte web sans privilège. Plutôt que d'exploiter le
serveur, il en lit la configuration, avec l'accord de M.~KOUNOU~: droits
sudo, programmes SUID, tâches planifiées, fichiers modifiables. Trois
vecteurs apparaissent~; un seul suffit.

\medskip
\textbf{Plan du chapitre.}
\begin{enumerate}
    \item Vertical et horizontal (section~\ref{sec:c11-types}).
    \item Les vecteurs sous Linux (section~\ref{sec:c11-linux}).
    \item Les vecteurs sous Windows (section~\ref{sec:c11-windows}).
    \item Lire un relevé (section~\ref{sec:c11-releve}).
    \item Corriger et détecter (section~\ref{sec:c11-corriger}).
\end{enumerate}

\separateur

% ============================================================================
\section{Vertical et Horizontal}
\label{sec:c11-types}
% ============================================================================

\begin{boxdef}
\termecle{Élévation de privilèges}~: obtention de droits supérieurs à ceux
du compte de départ (verticale~: utilisateur vers administrateur).
\termecle{Mouvement horizontal}~: passage d'un compte à un autre de même
niveau, souvent pour atteindre des données ou une autre machine.
\end{boxdef}

% ============================================================================
\section{Les Vecteurs sous Linux}
\label{sec:c11-linux}
% ============================================================================

\begin{tableaucours}{L{3.6cm}L{4.6cm}L{4.0cm}}{%
    \tableauHeader{Vecteur} &
    \tableauHeader{Faiblesse} &
    \tableauHeader{Correction}
}
sudo trop large & Commande autorisée sans mot de passe qui permet d'ouvrir
un interpréteur & Restreindre aux commandes nécessaires \\
Programme SUID inhabituel & S'exécute avec les droits de son propriétaire
& Retirer le bit SUID, supprimer le fichier \\
Tâche planifiée modifiable & Script exécuté par root, modifiable par
d'autres & Droits stricts sur les scripts \\
Capacités excessives & Un binaire dispose d'une capacité qui contourne les
contrôles & Retirer la capacité \\
Fichier sensible lisible & Clé, mot de passe dans un fichier accessible &
Droits stricts, coffre de secrets \\
Noyau non corrigé & Faille du noyau & Correctifs~; vecteur risqué à
exploiter \\
\end{tableaucours}
\vspace{-0.8em}
\begin{center}{\small\itshape Tableau~11.1. Vecteurs d'élévation sous Linux}\end{center}

% ============================================================================
\section{Les Vecteurs sous Windows}
\label{sec:c11-windows}
% ============================================================================

\begin{tableaucours}{L{3.6cm}L{4.6cm}L{4.0cm}}{%
    \tableauHeader{Vecteur} &
    \tableauHeader{Faiblesse} &
    \tableauHeader{Correction}
}
Service mal protégé & Binaire ou configuration du service modifiable par
un utilisateur & Droits stricts~; comptes de service dédiés \\
Chemin de service sans guillemets & Chemin avec espaces interprété par
morceaux & Mettre le chemin entre guillemets \\
Installation toujours élevée & Stratégie qui installe tout paquet avec les
droits système & Désactiver la stratégie \\
Mot de passe en clair & Script ou fichier de configuration lisible &
Supprimer, changer le mot de passe \\
Administrateur local partagé & Même mot de passe sur toutes les machines
& Mot de passe unique par machine \\
\end{tableaucours}
\vspace{-0.8em}
\begin{center}{\small\itshape Tableau~11.2. Vecteurs d'élévation sous Windows}\end{center}

\begin{boxscenario}{Opération LOGONE. Le script du partage}
\textbf{Constatation.} Sur le partage de fichiers \texttt{\textbackslash\textbackslash FS-01\textbackslash scripts},
lisible par tous les utilisateurs, un script de déploiement contient en
clair le mot de passe du compte administrateur local des postes de Douala.

\medskip
\textbf{Analyse} (confiance élevée). Depuis mars, chaque poste a son propre
mot de passe administrateur local (LAPS)~; mais ce script, antérieur, a
servi à leur installation, et 23 postes non réinstallés gardent l'ancien
mot de passe commun. C'est le constat critique T-02.
\end{boxscenario}

% ============================================================================
\section{Lire un Relevé}
\label{sec:c11-releve}
% ============================================================================

\begin{boxprocedure}
\textbf{Recenser les vecteurs à partir d'un relevé}
\begin{enumerate}
    \item Comptes et groupes~: qui a l'UID~0 (Linux), qui est administrateur
          local (Windows)~?
    \item Droits délégués~: \texttt{sudo -l}~; services et leurs droits.
    \item Programmes SUID et capacités hors des chemins du système.
    \item Tâches planifiées et scripts qu'elles exécutent~: qui peut les
          modifier~?
    \item Fichiers lisibles contenant des secrets.
    \item Version du noyau et correctifs, en dernier recours.
\end{enumerate}
\end{boxprocedure}

\begin{boxalert}
\textbf{Violation typique.} Commencer par une faille du noyau.

\smallskip
\textbf{Règle~:} une exploitation du noyau peut faire tomber le serveur. On
cherche d'abord les vecteurs de configuration, sans risque et plus
probables~; on rapporte la faille du noyau sans l'exploiter, sauf accord
explicite.
\end{boxalert}

% ============================================================================
\section{Corriger et Détecter}
\label{sec:c11-corriger}
% ============================================================================

\begin{boxobjet}{T-02 · Mot de passe administrateur local en clair}
\textbf{Description~:} un script lisible par tous sur
\texttt{\textbackslash\textbackslash FS-01\textbackslash scripts} contient le
mot de passe administrateur local commun de 23~postes de Douala.
\textbf{Preuve~:} chemin et date du fichier, mot de passe masqué.
\textbf{Gravité~:} critique. \textbf{Impact~:} administration de 23~postes
et mouvement latéral. \textbf{Recommandation~:} supprimer le script,
changer le mot de passe sur les 23~postes (LAPS), rechercher d'autres
secrets dans les partages. \textbf{Détection~:} connexions avec le compte
administrateur local depuis le réseau.
\end{boxobjet}

\begin{boxmetier}{Auditeur de sécurité}
\textbf{Chercher les secrets là où on travaille.} L'auditeur recherche les
mots de passe dans les partages, les scripts, les stratégies de groupe et
les dépôts internes~: c'est souvent là, et non dans les serveurs, que se
trouvent les clés du domaine.
\end{boxmetier}

\separateur

\begin{boxresume}
\textbf{Ce qu'il faut retenir de ce chapitre~:}
\begin{itemize}
    \item \textbf{Vertical} (plus de droits) et \textbf{horizontal} (autre
          compte).
    \item \textbf{Linux}~: sudo, SUID, tâches, capacités, fichiers, noyau.
    \item \textbf{Windows}~: services, chemins, stratégies, mots de passe en
          clair, administrateur local partagé.
    \item \textbf{On lit la configuration d'abord}~; le noyau en dernier.
    \item \textbf{Chaque vecteur} a sa correction et sa détection.
\end{itemize}
\end{boxresume}

\begin{boxinfo}
\textbf{Points de contrôle de l'élève.} Je sais~: distinguer élévation et
mouvement horizontal~; citer les vecteurs Linux et Windows~; lire un relevé
dans l'ordre~; rédiger un constat avec correction et détection.
\end{boxinfo}

\bigskip

\begin{boxexo}[title={\ding{45}\quad Exercice~11.1. Vertical ou horizontal \niveaubase}]
Classez~: passer du compte web au compte root~; utiliser le compte d'un
collègue de même service~; devenir administrateur local d'un poste.
\end{boxexo}

\begin{boxexo}[title={\ding{45}\quad Exercice~11.2. SUID \niveaubase}]
Pourquoi un fichier SUID appartenant à root dans \texttt{/tmp} est-il
suspect~? (cours SEC3266, lab de durcissement)
\end{boxexo}

\begin{boxexo}[title={\ding{45}\quad Exercice~11.3. Chemin \niveaubase}]
Expliquez le risque d'un chemin de service non entouré de guillemets.
\end{boxexo}

\begin{boxexo}[title={\ding{45}\quad Exercice~11.4. sudo \niveaumoyen}]
Un utilisateur peut lancer un éditeur de texte en sudo sans mot de passe.
Pourquoi est-ce dangereux~? Que proposez-vous~?
\end{boxexo}

\begin{boxexo}[title={\ding{45}\quad Exercice~11.5. Le script \niveaumoyen}]
Pourquoi 23~postes gardent-ils l'ancien mot de passe malgré LAPS~?
Comment les retrouver~?
\end{boxexo}

\begin{boxexo}[title={\ding{45}\quad Exercice~11.6. Ordre \niveaumoyen}]
Justifiez l'ordre de la procédure de recensement.
\end{boxexo}

\begin{boxexo}[title={\ding{45}\quad Exercice~11.7. Tâche planifiée \niveaumoyen}]
Une tâche exécutée par root lance un script modifiable par le groupe
« web ». Rédigez le constat.
\end{boxexo}

\begin{boxexo}[title={\ding{45}\quad Exercice~11.8. Détecter \niveaumoyen}]
Quels événements Windows permettent de détecter l'usage du compte
administrateur local depuis le réseau~?
\end{boxexo}

\begin{boxexo}[title={\ding{45}\quad Exercice~11.9. Recherche de secrets \niveauexpert}]
Concevez une campagne interne, autorisée, de recherche de secrets dans les
partages de LiZbiZ~: périmètre, méthode, traitement des résultats.
\end{boxexo}

\begin{boxexo}[title={\ding{45}\quad Exercice~11.10. Standard \niveauexpert}]
Rédigez cinq règles d'un standard de configuration qui ferment les vecteurs
des Tableaux~11.1 et~11.2.
\end{boxexo}

\bigskip

\begin{boxinfo}
\textbf{Pour aller plus loin.}
\begin{itemize}
    \item MITRE ATT\&CK~\cite{mitreattack}~; CIS Benchmarks~\cite{cis2021controls}.
\end{itemize}
\end{boxinfo}

\begin{boxinfo}
\textbf{Transition.} Une fois des droits obtenus, l'attaquant s'étend et
s'installe~: c'est la post-exploitation, objet du Chapitre~\ref{chap:post}.
\end{boxinfo}
